\section{Introduction}
Walking over irregular ground can cause a rigidly mounted camera to undergo periodic translation and rotation during swing--stance transitions, body weight transfer, and foot touchdowns. For an onboard imaging system, the consequences include unwanted inter-frame motion and image motion during exposure. These effects are relevant even when the robot successfully tracks its commanded velocity and its mean roll or pitch remains near zero.

One response is to exploit the leg configuration to stabilize the body. However, stance feet constrain the available movement, and the remaining joint workspace changes throughout the gait. The motions required for gaze direction, obstacle clearance, and velocity tracking may compete for the same joints. Under these conditions, camera stabilization is not an independently actuated task as it would be with a dedicated gimbal.

Existing work has shown that legged robots can adjust body pose for directed vision \cite{shaw2019}, that gait learning can improve image-based perception \cite{saggar2007}, and that hexapod mechanical and gait parameters can be optimized to reduce body oscillations \cite{burkus2022,yin2024}. Hence neither visual objectives nor gait adaptation nor their generic combination can be presented as novel. The unresolved engineering decision addressed here is more specific: \emph{when two feasible gaits generate different disturbances and leave different body-correction authority, which gait yields the smaller camera disturbance after the body controller has acted?}

A gait producing the smallest uncontrolled angular-velocity norm need not minimize residual image motion. A larger, slow oscillation may be easy to compensate if the support configuration provides sufficient workspace, whereas a smaller but fast disturbance or a disturbance occurring at a constrained contact phase may remain in the camera image. Conversely, a candidate with favorable raw motion may consume nearly all the joint motion needed for stabilization. This creates a coupling between gait selection and the actual body-correction mechanism.

We propose a controller that predicts gait-conditioned camera disturbance, approximates how much of that disturbance a fixed, project-owned body-pose compensator can cancel under support and actuator constraints, and selects a feasible gait based on the predicted remaining image motion. This compensator is a planned research component, not an existing executable in the current standalone runner. The local image interaction model supplies a task-dependent projection of camera motion; it is not used to claim that a rigid camera requires high-rate image feedback. RGB image measurements are instead retained for correction, calibration, and independent validation.

The manuscript is organized around two technical propositions to be tested: (i) candidate-specific feasible compensation changes the ranking of gaits compared with ranking by uncompensated body motion; and (ii) using this ranking online improves image motion under matched achieved travel speed, without materially reducing terrain traversal success. The experiments and statistical comparisons needed for these propositions are specified in Section~\ref{sec:experiments}, and no empirical superiority is assumed in advance.
